\section{Bounded presentations and relative logarithmic frames}
\label{sec:bounded-presentations}

We begin the proof of Theorem~\ref{thm:bounded-fibre}.  Throughout this
section and the next two, its dimension, singularity and reference-index
bounds are fixed.  A constant called uniform depends only on these bounds,
and on an additional numerical input when this is stated explicitly.  Put
\(n=\operatorname{trdeg}_K L>0\).

\subsection{Finite extensions and projective presentations}

We first explain precisely how bounded finite extensions of \(K\) can be
used.  Since \(L/K\) is regular, \(L\otimes_K K'\) is a field for a finite
extension \(K'/K\); we denote it by \(LK'\).  Let \(u'\) prolong a divisorial
valuation \(u\) of \(K\), with its scale chosen so that \(u'|_K=u\).
Every divisorial \(w\) of \(L\) over \(u\) has a compatible prolongation
to \(LK'\).  Indeed, tensor the completion at \(w\) with the completion
at \(u'\) over the completion at \(u\), choose a field factor, and use its
extended discrete valuation with the original scale.  Its restriction to
\(LK'\) has the required restrictions.  It is divisorial, and the
finite-extension discrepancy formula from Section~\ref{sec:tower} leaves
the discrepancies of pulled-back forms unchanged.  Consequently, for
each prolongation separately,
\begin{equation}
 \inf_{w'|_{K'}=u'} A_{\sigma_{LK'}}(w')
 =\inf_{w|_K=u}A_\sigma(w).
 \label{bounded:finite-extension}
\end{equation}
The two inequalities follow respectively by restriction and by the
compatible prolongation just constructed.  Prolonging any given valuation
also proves that homogeneous discrepancy inequalities descend, including
inequalities at valuations trivial on the base.

If \(u\) is integral-valued, its prolonged value group has denominator
dividing a ramification index at most \([K':K]\).  Multiplication by that
index restores integrality, at a bounded cost in a smallness threshold.
For the first assertion of Theorem~\ref{thm:bounded-fibre}, an upstairs
form descends by its norm: multiply its conjugate pluri-top forms in a
normal closure.  Dividing its divisor by the resulting tensor index gives
the average of the conjugate normalized divisors.  Formula
\eqref{bounded:finite-extension} therefore preserves both discrepancy
bounds, and the tensor index is multiplied by a bounded degree.  One may
equivalently use the usual norm over the given finite separable extension.

By BAB boundedness \cite[Theorem~1.1]{BAB}, geometric
\(\epsilon_0\)-lc Fano varieties of the relevant dimensions form a
bounded family. Hence \(G\) admits a very ample polarization of uniformly
bounded degree. It can be defined over a bounded extension of \(K\).
Here is a descent
argument that does not presume a polarization over \(K\).  The geometric
Picard group is finitely generated of bounded rank: klt Fano vanishing
gives \(H^1(G_{\overline K},\cO)=0\), and the exponential sequence, after
identifying the algebraically closed field with \(\C\), bounds its rank
by the Betti ranks in a bounded complex family.  It is torsion-free.  In
fact, if a line bundle \(T\) is torsion, polynomiality of Euler
characteristic gives \(\chi(T)=\chi(\cO)=1\), and klt Fano vanishing
applies to \(T\) as well.  Thus \(T\) has a nonzero section; an effective
numerically trivial divisor is zero, so \(T\) is trivial.

The continuous Galois action on this free group has finite image of
uniformly bounded order, since finite subgroups of \(\mathrm{GL}_q(\Z)\)
inject under reduction modulo \(3\).  After an extension killing that
action, the polarization class is invariant.  Its complete linear system
descends to a possibly twisted projective space: the isomorphisms of the
polarization differ only by scalars, so their projectivizations satisfy
descent.  The dimension of this Severi--Brauer variety is bounded by the
degree and dimension of \(G\).  A splitting extension of degree at most
the degree of its central simple algebra is therefore bounded.  We may
and do use an embedding over the resulting field, with
\(A=\cO_G(1)\), whose ambient dimension and degree are bounded.

Choose a nonzero rational absolute top differential \(\zeta\) on \(K\).
The differential determinant for \(L/K/\C\) writes the reference form as
\begin{equation}
 \sigma=\eta\otimes\zeta^{\otimes r_0},\qquad 1\le r_0\le r,
 \label{bounded:relative-form}
\end{equation}
where \(\eta\) is a rational relative \(r_0\)-pluri-top form.
Its divisor on \(G\) is nonpositive, and its total absolute degree is
uniformly bounded.  Indeed,
\[
 \deg_A\!\left(-\frac1{r_0}\operatorname{div}_G\eta\right)
 =-K_G\cdot A^{n-1};
\]
adjunction to a general complete-intersection curve in the smooth locus
bounds this intersection in the bounded embedded family.  For \(n=1\)
the normal curve itself is smooth and the same formula applies.

Take a finite linear projection with affine coordinates
\(z^0_1,\ldots,z^0_n\).  The differential
\(dz^0_1\wedge\cdots\wedge dz^0_n\) has poles bounded by a fixed
multiple of the infinity hyperplane, as is seen by differentiating at
each codimension-one point.  Its zeros have bounded degree as well, by
the preceding canonical-degree calculation.  Hence
\begin{equation}
 \eta=(dz^0_1\wedge\cdots\wedge dz^0_n)^{\otimes r_0}\frac{P}{Q},
 \label{bounded:polynomial-form}
\end{equation}
where \(P,Q\) are restrictions of homogeneous polynomials of the same
uniformly bounded degree.

For completeness, a rational function with pole divisor of degree at
most \(a\) on one of these normal embedded varieties is a quotient of
polynomial restrictions of degree at most \(C(1+a)\).  Cut each pole
prime by a hypersurface containing it but not the variety, of degree at
most the degree of the prime; a general finite linear projection gives
such a hypersurface.  Multiply the equations with the pole
multiplicities.  Normality makes the resulting numerator a regular
section of the corresponding twist.  Uniform regularity in the bounded
projective Hilbert families lifts it to an ambient polynomial after a
uniform further twist.  Applying this argument to \eqref{bounded:relative-form}
gives \eqref{bounded:polynomial-form}.  Applying it to a ratio of relative
forms gives a second bound which will be used repeatedly: if another
form has effective boundary on \(G\), and the forms are compared at a
common tensor index \(p\), then their ratio has a polynomial presentation
of degree at most \(Cp\), independently of the index of that other form.

\subsection{Preparation valid at every parameter}

Hilbert schemes and the coefficients in \eqref{bounded:polynomial-form}
give finitely many algebraic families of varieties with rational
relative forms.  We retain the scalar coefficient of each form, so the
form itself, and not just its divisor, specializes correctly.  Their
parameter spaces may be taken to be dense opens of projective varieties
over \(\C\).  We will repeatedly compactify dominantly, resolve, and
shrink these opens.

These operations can be arranged to cover every field-valued parameter
under consideration.  To see this, take reduced closures of the valid
parameter points.  On a dense open of each irreducible component,
flatness and the open geometric loci give normal, geometrically integral
fibres.  Resolve the generic fibre together with the relative form
divisor and spread this resolution out.  After shrinking, its relative
divisor has fixed simple-normal-crossings coefficients and specializes
as the divisor of the fibre form; the exceptional and divisorial images
are also constant in the required sense.  Effectivity of the boundary
and log canonicity are thus tested on an open.  Density of the valid
points forces these properties at the generic point, so this open
contains all the desired properties.  This argument uses no bounded
index for the underlying \(K_G\): the adjoint with the form boundary is
already \(\Q\)-principal.

We also spread all inverse rational maps, their identities on dense
opens, the exact differential and form identities, and their nonzero
denominators.  Shrinking makes these identities valid for every
field-valued point of the open, including a point whose map to the
parameter space is not dominant.  The finitely many irreducible closed
complements have smaller dimension.  Noetherian induction applied to
those complements produces finitely many such packages.  In particular,
any later shrinking is accompanied by this induction; no parameter is
discarded.

\begin{proposition}[Relative logarithmic preparation]
\label{bounded:preparation}
The finitely many packages can be chosen with projective morphisms
\(f:Y\to B\) having the following properties.  The base is smooth and
strict toroidal; the source is strict simplicial toroidal; and every
source cone maps onto a base cone.  The good parameter open is disjoint
from the base boundary.  For each field-valued point used in a package,
\(Y_K\) maps projectively and birationally to the prescribed \(G\), and
the differential and form identities specialize.  The divisor of
\(\eta\), regarded as a relative logarithmic pluri-top form, is
toroidal.  Arbitrary prescribed proper closed subsets can be included
in the boundary during this construction.
\end{proposition}

\begin{proof}
We use weak toroidalization \cite[Theorem~1.1]{ADK} for projective dominant morphisms
in characteristic zero, with prescribed closed subsets, and prove the
additional assertion about the logarithmic form.  Only projective
birational modifications of the source over the resulting base change
are needed.

\smallskip\noindent\emph{Induction and the relative form.}
In relative dimension zero, mark the zeros and poles of the relative
form function.  A dominant generically finite toroidal morphism has
invertible logarithmic differential determinant in characteristic zero.
This follows in toric charts from the full rank of the map of torus
lattices, including the unit directions.  More generally that map has
full target rank for a dominant toroidal morphism: otherwise the map on
logarithmic differentials in the full toric frames would be generically
rank-deficient even after completion, contrary to dominance and
separability.  Faithful flatness of completion transfers this rank test
to the original chart.

For the inductive step, choose independent relative rational functions
and resolve their graph to factor the morphism projectively through an
intermediate base of relative dimension one less.  The top arrow then
has relative dimension one.  In the determinant factorization write
\(\eta=\eta_1\otimes\eta_{\mathrm{low}}\), with the appropriate common
tensor power understood.  Apply weak toroidalization to the top arrow,
marking the prescribed sets and the closures of the support of the
divisor of \(\eta_1\) on its original normal proper generic curve.
A birational normal proper model of that generic curve is the same
curve.  Consequently every component of the relative logarithmic
divisor of \(\eta_1\) which is not already in the boundary is vertical.

Refine this top toroidal arrow so that its base \(B_1\) is smooth
strict toroidal, its source is strict simplicial toroidal, and all cone
maps are onto cones.  Mark on \(B_1\) the images of those finitely many
bad vertical components and the entire boundary of \(B_1\).  Apply the
induction hypothesis to the lower arrow \(B_1\to B_{\mathrm{old}}\).
It gives an arrow \(B_2\to B_{\mathrm{new}}\), mapping to the old lower
arrow, with the logarithmic form property for
\(\eta_{\mathrm{low}}\).

\smallskip\noindent\emph{Normalized main pullback.}
We verify the normalized main pullback of the top arrow under
\(B_2\to B_1\).  This map need not be toroidal, or dominant on every
subsequent parameter.  The marked inverse images nevertheless ensure
that the old boundary equations pull back to monomials times units in
complete toroidal charts.  One can see this using the log structure of
functions invertible off the boundary.  Alternatively, on a strict
simplicial chart some multiple of the pull divisor is the divisor of a
character monomial.  Normality makes the corresponding power of the
function that monomial times a unit.  At a geometric closed point the
unit has a formal root in characteristic zero.  The first term for
generic positive character weights shows that the monomial exponent
itself is divisible by that power, proving the assertion for the
original function.

At such a point of the old top arrow, write the boundary coordinates of
the smooth base as monomials \(z^{p_j}\) times units in the source
monoid.  The \(p_j\) are linearly independent.  The smooth base
parameters transverse to its boundary can be made part of the smooth
source parameter system: the stratum map is smooth, as follows from
the full-rank torus map in characteristic zero.  First absorb the old
source units by a unit-valued character on the source lattice, taking
formal roots when its finite lattice index requires them.  Scaling
monomials by such units is an automorphism of the complete local ring;
its action on the monoid generators in the cotangent space is
invertible, and completeness gives an inverse by successive
approximation.  Combined with the preceding smooth-parameter change,
the cotangent matrix is invertible.  The logarithmic Jacobian of this
old source-coordinate change is a unit.  Its relative differential
need not be zero, but it changes a logarithmic determinant frame by a
unit and hence does not change its divisor.

Now pull back to \(B_2\).  Absorb the remaining unit factors, which come
from the new base, along the old source characters.  These multipliers
have zero relative differential over \(B_2\).  Substitute for the old
smooth base parameters, retaining all the smooth parameters of
\(B_2\).  The remaining equations are binomial:
\begin{equation}
 z^{p_j}=(z')^{q_j}.
 \label{bounded:binomial-pullback}
\end{equation}
Their normalized main branches are toroidal, with cone the intersection
of the product cones with
\[
 p_*x=q_*y.
\]
This intersection meets the relative interior of the product cone:
the first cone maps onto the old base cone, and the relative interior
of the second maps into the relative interior of its relevant old
base cone.  Inverting all monomials in
\eqref{bounded:binomial-pullback} gives the character-group pushout.
Its finite torsion, when present, separates the torus branches.  The
closure of each such branch is an affine monoid variety; its
normalization is the saturation of that monoid in its character
lattice.  The interior condition makes its monoid sharp and gives a
unique point over the monoid origin on each branch.  Completion at
that point is one normal monoid-series domain, with the free smooth
parameters adjoined.

These are the branches of the actual algebraic main pullback.  The
joint nonboundary open is dense in that main component, so boundary
monomials are nonzerodivisors there.  Flat completion preserves this
property; no minimal completed branch can be supported solely on
vanishing boundary monomials.  Thus the relevant completed branches
are precisely the closures of the torus branches just computed.
Excellence gives finite normalization and its compatibility with this
completion; geometric generic integrality identifies the prescribed
main generic component.  These facts also justify performing the
binomial computation first in polynomial monoid algebras and then
completing.  They do not require an arbitrary formal coordinate change
to be algebraic.

The chart map to \(B_2\) retains its given character map and smooth
parameters.  Its cone is onto the \(B_2\) cone, since the old top cone
was onto.  Characters rationally complementary to the \(p_j\), together
with the remaining relative smooth parameters, give the same relative
logarithmic determinant frame after pullback: saturation changes no
rational span, and the new unit multipliers contribute only base
differentials.  It follows that the relative logarithmic determinant
pulls back unchanged.  The divisor of \(\eta_1\) is now wholly
supported in the boundary, because its bad vertical images were
marked.  Composing with \(B_2\to B_{\mathrm{new}}\) gives a toroidal
map; in these same charts its base equations are monomials, its log
rank is full, and its stratum coordinates are smooth.  The relative
determinant factorization proves the assertion for \(\eta\).  The
models continue to map birationally to the old ones and retain all
inverse marks.

\smallskip\noindent\emph{Projective cone refinement.}
We carry out these refinements on the finite-type packages over
\(\C\), before evaluating at field-valued parameter points. First take
compatible projective toroidal resolutions, so that the source and
base are smooth and strict. On the resulting regular source cone
complex, assign value \(1\) to each primitive ray generator and extend
linearly on each cone. These functions agree on common faces and
define a slicing function \(h\), positive away from the origins.
Projective equidimensionalization
\cite[Proposition~4.4 and Lemma~4.1]{AbramovichKaruRefinement}
then gives a smooth base and makes every source cone map onto a base
cone. Retain \(h\) on the source subdivision: every new cone lies in
an old cone, so \(h\) remains linear there and positive away from the
origins. Apply
\cite[Corollary~0.11]{AbramovichRojasTriangulations}
to its \(1\)-skeleton with the unchanged triangulation induced by the
zero lifting function. This gives a projective triangulation with
no new rays. Since the base is simplicial, each old source ray maps
to a base ray or to zero
\cite[Remark~4.2]{AbramovichKaruRefinement}; the image of each new
simplicial cone is therefore a face of a base cone. The final source
is strict and simplicial, and the cone-surjectivity condition is
preserved. All these are toroidal modifications, so the inverse marks
and the pullback identities for the logarithmic divisors are retained.

\smallskip\noindent\emph{Return to all parameters.}
Apply the construction to the compactified families, and use the
all-parameter shrinking and induction above. Any loci newly excluded
by these refinements undergo the same noetherian restart before the
finite package list and its common moduli multiple are fixed. On the good
isomorphism open the parameter map already lifts to the modified
base; properness then extends a valuation-ring map when required.
The retained inverse maps and form identities give the claimed
birational and differential assertions on every fibre used.
\end{proof}

All models, cones and chart functions in Proposition~\ref{bounded:preparation}
belong to finitely many finite-type families.  After the indicated
shrinking and spreading, the generic models \(Y_K\) are bounded in
fixed projective embeddings.  Pullback of the polynomial presentations
on \(G\) shows that a comparison function at tensor index \(p\) has
generic zero and pole divisors on \(Y_K\) of combined degree at most
\(Cp\), counted with multiplicities.

\subsection{The minimum function}

Let \(\phi\) denote the toroidal order function of
\(r_0^{-1}\operatorname{div}^{\log}_{Y/B}(\eta)\).  It is rational linear
on each source cone.  On every horizontal ray, that is, a ray mapping
to zero, it is nonnegative by log canonicity of the generic form pair.
For a base weight \(b\), define
\begin{equation}
 \mu(b)=\min_{f_*x=b}\phi(x),
 \label{bounded:mu}
\end{equation}
using all source cones over its base cone.

Every cone mapping onto the cone at a given base point can be used
over that point, with an appropriate ordinary-fibre component label.
Indeed, its smooth stratum map has open dense image in the base
stratum.  Properness makes its closure meet every point of that
stratum.  At a deeper point of this closure, the face corresponding
to the original stratum still maps onto the base cone.  In the formal
orbit-closure chart for that face, all base boundary monomials vanish
identically, and the remaining base parameters are smooth stratum
parameters which can be eliminated. The full-rank Jacobian in these
unconstrained stratum coordinates survives passage to the face
orbit-closure quotient. The open face stratum is dense
in this fibre chart.  Thus that fibre has actual points of the desired
open source stratum.  This last calculation, in addition to
properness, supplies the assertion at every point.  Points and
components may be taken geometrically.

The minimum in \eqref{bounded:mu} exists and is rational piecewise
linear.  On any simplicial cone, allocate each base coordinate to the
source rays above its axis and choose a ray of least cost.  Horizontal
rays have nonnegative cost and can be set to zero.  There are only
finitely many cones and choices, so taking their minimum gives the
claim.  The resulting minima are continuous across base faces.  A
candidate on a face can be approached from every incident deeper base
cone: properness specializes its source-stratum closure to a stratum
over that deeper cone, and arbitrarily small added weights in the
new directions supply the missing base coordinates.  Conversely, a
limit of minimum candidates from the interior supplies a face
candidate; the horizontal nonnegative directions can be omitted.
After compatible projective subdivisions, with the refinements in
Proposition~\ref{bounded:preparation}, we may therefore assume that
\(\mu\) is linear on each base cone.  Write
\begin{equation}
 M_0\text{ for the toroidal \(\Q\)-divisor on \(B\) with order function
 }\mu.
 \label{bounded:moduli}
\end{equation}

The next section compares this geometric minimum with actual
divisorial prolongations of a base valuation.  That comparison will
identify \(M_0\) as the moduli divisor and give the bounded-index
base form.
